\noindent \textbf{Two-stage TAL}. These approaches first generate candidate video segments as action proposals, and further classify the proposals into actions and refine their temporal boundaries. Previous works focused on action proposal generation, by either classifying anchor windows~\cite{caba2016fast,escorcia2016daps,sstcvpr2017} or detecting action boundaries~\cite{bsneccv2018,bmniccv2019,mggcvpr2019,gong2020scale,bottumuptaleccv2020}, and more recently using a graph representation~\cite{bcgnneccv2020,gtadcvpr2020} or Transformers~\cite{tan2021relaxed,chang2021augmented,wang2021temporal}. Others have integrated proposal generation and classification into a single model~\cite{shou2016temporal,ssniccv2017,cdc2017cvpr,fasterrcnncvpr2018}. More recent effort investigates the modeling of temporal context among proposals using graph neural networks~\cite{pgcniccv2019,gtadcvpr2020,vsgniccv2021} or attention and self-attention mechanisms~\cite{contextlociccv2021,tcanetcvpr2021,csaiccv2021}. Similar to previous approaches, our method considers the modeling of long-term temporal context, yet uses a self-attention within a Transformer model. Different from previous approaches, our model detects actions without using proposals.\smallskip

\noindent \textbf{Single-stage TAL}. Several recent works focused on single-stage TAL, seeking to localize actions in a single shot without using action proposals. Many of them are anchor-based (\eg using anchor windows sampled from sliding windows). Lin \et\ \cite{sstadmm2017} presented the first single-stage TAL using convolutional networks, borrowing ideas from a single-stage object detector~\cite{liu2016ssd}. Buch \et\ \cite{sstadbmvc2017} presented a recurrent memory module for single-stage TAL. Long \et\  \cite{gaussiancvpr2019} proposed to use Gaussian kernels to dynamically optimize the scale of each anchor, based on a 1D convolutional network. Yang \et\ \cite{a2nettip2020} explored the combination of anchor-based and anchor-free models for single-stage TAL, again using convolutional networks. More recently, Lin \et\ \cite{afsdcvpr2021} proposed an anchor-free single-stage model by designing a saliency-based refinement module incorporated in convolutional network. Similar ideas were also explored in video grounding~\cite{zeng2020dense}. 

Our model falls into the category of single-stage TAL. Indeed, our formulation follows a minimalist design of sequence labeling by classifying every moment and regressing their action boundaries, previously discussed in~\cite{a2nettip2020,afsdcvpr2021}. The key difference is that we design a Transformer network for action localization. The result is a single stage anchor-free model that outperforms all previous methods. A concurrent work from Liu \et\ \cite{tadtrarxiv2021} also used Transformer for TAL, yet considered a set prediction problem similar to DETR~\cite{detreccv2020}. %
\smallskip 


\noindent \textbf{Spatial-temporal Action Localization}.
A related yet different task, known as spatial-temporal action localization, is to detect the actions both temporally and spatially, in the form of moving bounding boxes of an actor. It is possible that TAL might be used as a first step for spatial-temporal localization. Girdhar \et\ \cite{actiontransformercvpr2019} proposed to use Transformer for spatial-temporal action localization. While both our work and~\cite{actiontransformercvpr2019} use Transformer, the two models differ significantly. We consider a sequence of video frames as the inputs, while~\cite{actiontransformercvpr2019} used a set of 2D object proposals. Moreover, our work addresses a different task of TAL.\smallskip

\noindent \textbf{Object Detection}. TAL models have been heavily influenced by the developments of object detection models. Some of our model design, including the multiscale feature representation and convolutional decoder, is inspired by feature pyramid network~\cite{fpncvpr2017} and RetinaNet~\cite{focaliccv2017}. Our training using center sampling also stems from recent single-stage object detectors~\cite{duan2019centernet,fcosiccv2019,zhang2020bridging}.\smallskip

\noindent \textbf{Vision Transformer}. Transformer models were originally developed for NLP tasks~\cite{transformernips2017}, and has demonstrated recent success for many vision tasks. ViT~\cite{vitarxiv2020} presented the first pure Transformer-based model that can achieve state-of-the-art performances on image classification. Subsequent works, including DeiT~\cite{deiticml2021}, T2T-ViT~\cite{t2tviticcv2021}, Swin Transformer~\cite{swiniccv2021}, Focal Transformer~\cite{yang2021focal} and PVT~\cite{pvticcv2021}, have further pushed the envelope, resulting in vision Transformer backbones with impressive results on classification, segmentation, and detection tasks. Transformer have also been explored in object detection~\cite{detreccv2020,deformabledetriclr2021,dynamicdetriccv2021,pnpiccv2021}, semantic segmentation~\cite{transsegcvpr2021,xie2021segformer,cheng2021per}, and video representation learning~\cite{arnab2021vivit,liu2021video,fan2021multiscale}. Our model builds on these developments and presents one of the first Transformer models for TAL.  